% macros.tex — theorem environments (aliascnt-safe) + notation macros.
% No \documentclass / \begin{document} here; this is \input{}-ed by main.tex.

% ---------------------------------------------------------------------------
% Theorem environments — aliascnt so cleveref keys names on the right counter.
% ---------------------------------------------------------------------------
\theoremstyle{plain}
\newtheorem{theorem}{Theorem}[section]

\newaliascnt{lemma}{theorem}
\newtheorem{lemma}[lemma]{Lemma}
\aliascntresetthe{lemma}

\newaliascnt{proposition}{theorem}
\newtheorem{proposition}[proposition]{Proposition}
\aliascntresetthe{proposition}

\newaliascnt{corollary}{theorem}
\newtheorem{corollary}[corollary]{Corollary}
\aliascntresetthe{corollary}

\theoremstyle{definition}
\newaliascnt{definition}{theorem}
\newtheorem{definition}[definition]{Definition}
\aliascntresetthe{definition}

\newaliascnt{assumption}{theorem}
\newtheorem{assumption}[assumption]{Assumption}
\aliascntresetthe{assumption}

\newaliascnt{fact}{theorem}
\newtheorem{fact}[fact]{Fact}
\aliascntresetthe{fact}

\theoremstyle{remark}
\newaliascnt{remark}{theorem}
\newtheorem{remark}[remark]{Remark}
\aliascntresetthe{remark}

% Cleveref names — keyed on the aliascnt counters above.
\crefname{theorem}{Theorem}{Theorems}
\crefname{lemma}{Lemma}{Lemmas}
\crefname{proposition}{Proposition}{Propositions}
\crefname{corollary}{Corollary}{Corollaries}
\crefname{definition}{Definition}{Definitions}
\crefname{assumption}{Assumption}{Assumptions}
\crefname{fact}{Fact}{Facts}
\crefname{remark}{Remark}{Remarks}
\crefname{equation}{Eq.}{Eqs.}

% ---------------------------------------------------------------------------
% Asymptotics
% ---------------------------------------------------------------------------
\newcommand{\Otil}{\widetilde{O}}

% ---------------------------------------------------------------------------
% Probability / expectation
% ---------------------------------------------------------------------------
\newcommand{\E}{\mathbb{E}}
\newcommand{\Prob}{\mathbb{P}}
\newcommand{\R}{\mathbb{R}}

% ---------------------------------------------------------------------------
% Norms / inner products
% ---------------------------------------------------------------------------
\newcommand{\norm}[1]{\left\| #1 \right\|}
\newcommand{\Norm}[1]{\left\| #1 \right\|}        % alias for big norm in displays
\newcommand{\inner}[2]{\left\langle #1, #2 \right\rangle}
\newcommand{\Inner}[2]{\left\langle #1,\, #2 \right\rangle} % big inner product
\newcommand{\Mnorm}[2]{\left\| #1 \right\|_{#2}}  % Mahalanobis / weighted norm
\DeclareMathOperator*{\argmax}{arg\,max}

% ---------------------------------------------------------------------------
% MDP / RL objects
% ---------------------------------------------------------------------------
\newcommand{\St}{\mathcal{S}}            % state space
\newcommand{\Ac}{\mathcal{A}}            % action space
\newcommand{\featmap}{\phi}              % feature map
\newcommand{\Gram}{\Lambda}              % Gram / covariance matrix
\newcommand{\Ev}{\mathcal{E}}            % successful event
\newcommand{\Vsp}{\mathcal{V}}           % value-function class
\newcommand{\Pop}{P}                     % transition kernel
\newcommand{\Reg}{\mathrm{Regret}}
\newcommand{\T}{\mathbb{T}}              % Bellman backup operator
